%%%%%%%%%%%%%%%%%%%%%%%%%%%%%%%%%%%%%%%%%%%%%%%%%%%%%%%%%%%%%%%%%%%%%%%%%%%%%%%
% Full instructions available at:
% https://github.com/elauksap/focus-beamertheme

\documentclass{beamer}
\usetheme{focus}

% package for font
\usepackage{fontspec}
\defaultfontfeatures{Mapping=tex-text}  %%如果没有它，会有一些 tex 特殊字符无法正常使用，比如连字符。
\usepackage{xunicode,xltxtra}
\usepackage[BoldFont,SlantFont,CJKnumber,CJKchecksingle]{xeCJK}  % \CJKnumber{12345}: 一万二千三百四十五
\usepackage{CJKfntef}  %%实现对汉字加点、下划线等。
\usepackage{pifont}  % \ding{}
% package for math
\usepackage{amsfonts}
% package for graphics
\usepackage[americaninductors,europeanresistors]{circuitikz}
\usepackage{tikz}
\usetikzlibrary{plotmarks}  % placements=positioning
\usepackage{graphicx}  % \includegraphics[]{}
\usepackage{subfigure}  %%图形或表格并排排列
% package for table
\usepackage{colortbl,dcolumn}  %% 彩色表格
\usepackage{multirow}
\usepackage{multicol}
\usepackage{booktabs}
% package for code
\usepackage{fancyvrb}
\usepackage{listings}
% setting for font
%%\setCJKmainfont{Adobe Kaiti Std}
\setCJKmainfont{SimSun}
%%\setCJKmainfont{Microsoft YaHei}
%%\setCJKmainfont{FangSong_GB2312} 
%%\setmainfont{Times New Roman}
%\setsansfont[Mapping=tex-text]{Adobe Song Std}
%如果装了Adobe Acrobat，可在font.conf中配置Adobe字体的路径以使用其中文字体。
%也可直接使用系统中的中文字体如SimSun、SimHei、微软雅黑等。
%原来beamer用的字体是sans family；注意Mapping的大小写，不能写错。
%设置字体时也可以直接用字体名，以下三种方式等同：
%\setromanfont[BoldFont={黑体}]{宋体}
%\setromanfont[BoldFont={SimHei}]{SimSun}
%\setromanfont[BoldFont={"[simhei.ttf]"}]{"[simsun.ttc]"}

% font setting by xeCJK
\setCJKfamilyfont{NSimSun}{NSimSun}
\newcommand{\song}{\CJKfamily{NSimSun}}
%%%\setCJKfamilyfont{AdobeSongStd}{Adobe Song Std}
%%%\newcommand{\AdobeSong}{\CJKfamily{AdobeSongStd}}
\setCJKfamilyfont{FangSong}{FangSong_GB2312}
\newcommand{\fang}{\CJKfamily{FangSong}}
%%%\setCJKfamilyfont{AdobeFangsongStd}{Adobe Fangsong Std}
%%%\newcommand{\AdobeFang}{\CJKfamily{AdobeFangsongStd}}
\setCJKfamilyfont{SimHei}{SimHei}
\newcommand{\hei}{\CJKfamily{SimHei}}
%%%\setCJKfamilyfont{AdobeHeitiStd}{Adobe Heiti Std}
%%%\newcommand{\AdobeHei}{\CJKfamily{AdobeHeitiStd}}
\setCJKfamilyfont{KaiTi}{KaiTi}
\newcommand{\kai}{\CJKfamily{KaiTi}}
%%%\setCJKfamilyfont{AdobeKaitiStd}{Adobe Kaiti Std}
\newcommand{\AdobeKai}{\CJKfamily{AdobeKaitiStd}}
\setCJKfamilyfont{LiSu}{LiSu}
\newcommand{\li}{\CJKfamily{LiSu}}
\setCJKfamilyfont{YouYuan}{YouYuan}
\newcommand{\you}{\CJKfamily{YouYuan}}
\setCJKfamilyfont{FZJingLei}{FZJingLeiS-R-GB}
\newcommand{\jinglei}{\CJKfamily{FZJingLei}}
\setCJKfamilyfont{MSYH}{Microsoft YaHei}
\newcommand{\msyh}{\CJKfamily{MSYH}}

\setbeamerfont{frametitle}{series=\Large\hei} % 修改Beamer标题字体

\graphicspath{{figures/}}  %%图片路径
\renewcommand\figurename{图}

% setting for pdf
\hypersetup{% pdfpagemode=FullScreen,%
  pdfauthor={Xiaodong Wang},%
  pdftitle={Title},%
  CJKbookmarks=true,%
  bookmarksnumbered=true,%
  bookmarksopen=false,%
  plainpages=false,%
  colorlinks=true,%
  citecolor=green,%
  filecolor=magenta,%
  linkcolor=blue,%red(default)
  urlcolor=cyan}

% setting for fontspec
\XeTeXlinebreaklocale "zh"  %%表示用中文的断行
\XeTeXlinebreakskip = 0pt plus 1pt minus 0.1pt  %%多一点调整的空间
%%%%%

% 自定义颜色
\def\Red{\color{red}}
\def\Green{\color{green}}
\def\Blue{\color{blue}}
\def\Mage{\color{magenta}}
\def\Cyan{\color{cyan}}
\def\Brown{\color{brown}}
\def\White{\color{white}}
\def\Black{\color{black}}

\lstnewenvironment{xmlCode}[1][]{% for Java
  \lstset{
    basicstyle=\tiny\ttfamily,%
    columns=flexible,%
    framexleftmargin=.7mm, %
    % frame=shadowbox,%
    % rulesepcolor=\color{cyan},%
     frame=single,%
    backgroundcolor=\color{white},%
    xleftmargin=4\fboxsep,%
    xrightmargin=4\fboxsep,%
    numbers=left,numberstyle=\tiny,%
    numberblanklines=false,numbersep=7pt,%
    language=xml, %
    }\lstset{#1}}{}

\lstnewenvironment{javaCode}[1][]{% for Java
  \lstset{
    basicstyle=\tiny\ttfamily,%
    columns=flexible,%
    framexleftmargin=.7mm, %
    frame=shadowbox,%
    rulesepcolor=\color{cyan},%
    % frame=single,%
    backgroundcolor=\color{white},%
    xleftmargin=4\fboxsep,%
    xrightmargin=4\fboxsep,%
    numbers=left,numberstyle=\tiny,%
    numberblanklines=false,numbersep=7pt,%
    language=Java, %
    }\lstset{#1}}{}

\lstnewenvironment{shCode}[1][]{% for Java
  \lstset{
    basicstyle=\scriptsize\ttfamily,%
    columns=flexible,%
    framexleftmargin=.7mm, %
    frame=shadowbox,%
    rulesepcolor=\color{brown},%
    % frame=single,%
    backgroundcolor=\color{white},%
    xleftmargin=4\fboxsep,%
    xrightmargin=4\fboxsep,%
    numbers=left,numberstyle=\tiny,%
    numberblanklines=false,numbersep=7pt,%
    language=sh, %
    }\lstset{#1}}{}

\newcommand\ask[1]{\vskip 4bp \tikz \node[rectangle,rounded corners,minimum size=6mm,
  fill=white,]{\Cyan \includegraphics[height=1.5cm]{question} \Large \msyh #1};}

\newcommand\wxd[1]{\vskip 4bp \tikz \node[rectangle,minimum size=6mm,
  fill=blue!60!white,]{\White \ding{118} \msyh #1};}

\newcommand\cxf[1]{\vskip 4bp \tikz \node[rectangle,rounded corners,minimum size=6mm,
  fill=purple!60!white,]{\White \ding{42} \msyh #1};}

\newcommand\xyy[1]{\vskip 2bp \tikz \node[rectangle,minimum size=3mm,
  fill=black!80!white,]{\White \msyh\scriptsize #1};}

\newcommand\tta[1]{\vskip 4bp \tikz \node[rectangle,minimum size=6mm,
  fill=blue!60!white,]{\White \ding{118} \msyh #1};}

\newcommand\ttb[1]{\vskip 4bp \tikz \node[rectangle,rounded corners,minimum size=6mm,
  fill=purple!60!white,]{\White \ding{42} \msyh #1};}

\newcommand\ttc[1]{\vskip 2bp \tikz \node[rectangle,minimum size=3mm,
  fill=black!80!white,]{\White \msyh\scriptsize #1};}

\newcommand\homework[1]{\vskip 2bp \tikz \node[rectangle,minimum size=3mm,
  fill=red!80!white,]{\White \ding{45} \msyh\scriptsize 课后小作业 } ; {\kai\small #1}} 

\newcommand\notice[1]{\vskip 4bp \tikz \node[rectangle,rounded corners,minimum size=6mm,
  fill=red!80!white,]{\White \scriptsize \ding{42} \msyh #1};}

\newcommand\samp[1]{\vskip 2bp \tikz \node[rectangle,minimum size=3mm,]{\Mage\msyh \small CODE \ding{231} \Black #1};\vskip -8bp}

\newcommand\codeset[1]{\vskip 2bp \tikz \node[rectangle,minimum
  size=3mm]{\Mage\msyh \small 课程配套代码 \ding{231} \Black
    #1};\vskip -4bp}

\newcommand\pptlink[2]{\vskip 4bp \tikz \node[rectangle,rounded corners,minimum size=6mm,
  fill=blue!70!white,]{\href{run:#1}{\White \scriptsize \msyh 动画演示 #2}};}


\makeatletter
\newcommand{\Extend}[5]{\ext@arrow 0099{\arrowfill@#1#2#3}{#4}{#5}}
\makeatother
%%%%%%%%%%%%%%%%%%%%%%%%%%%%%%%%%%%%%%%%%%%%%%%%%%%%%%%%%%%%%%%%%%%%%%%%%%%%%%%
\title{\hei JAVA应用与开发\\  泛\quad 型}
\subtitle{\it 让我们愉快的Coding起来吧...}
\author{王晓东}
\titlegraphic{\vspace{-8em}\includegraphics[height=10cm]{static/ouc.pdf}\vspace{-6em}} 
\institute{中国海洋大学信息学院计算机系}
\date{\today}

\begin{document}

\begin{frame}
  \maketitle
\end{frame}

\begin{frame}
  \frametitle{学习目标}

  \begin{itemize}
  \item 理解泛型的概念，掌握其基本应用
    \begin{itemize}
    \item 集合框架中的泛型
    \item 泛型的向后兼容性
    \end{itemize}
  \item 掌握自定义泛型类和泛型方法
    \begin{itemize}
    \item 理解类型参数
    \item 理解差异性并能够定义自己的泛型类和泛型方法
    \item 受限制的类型参数
    \end{itemize}
  \item 学会处理泛型类型，包括使用通配符实现泛型容器遍历和操作
  \end{itemize}
\end{frame}

\begin{frame}
  \frametitle{大纲}
  \tableofcontents
\end{frame}

\section{泛型概念}

\begin{frame}[fragile] % [fragile]参数使得能够插入代码
  \frametitle{什么是泛型}

  \begin{block}{泛型（Generics）}
    泛型机制自JDK 5.0开始引入，其实质是{\hei\Blue 将原本确定不变的数据
      类型参数化}。

    作为对原有Java类型体系的扩充，使用泛型可以提高Java应用程序的类型安
    全、可维护性和可靠性。
  \end{block}
\end{frame}

\begin{frame}[fragile] 
  \frametitle{什么是泛型}

  \tta{集合框架中的数据造型问题}

  传统的集合容器为了提供广泛的适用性，会将所有加入其中的元素当
  作Object类型来处理。基于此原因，在实际使用时，我们必须将从集合中取出
  的元素值再强制转换（造型）为所期望的类型。

  \pause
  
  \notice{无泛型机制的集合容器}
  
  \begin{javaCode}
    Vector v = new Vector();
    v.addElement(new Person("Tom", 18));
    Person p = (Person) v.elementAt(0);
    p.showInfo();
  \end{javaCode}
  
\end{frame}

\begin{frame}[fragile]
  \frametitle{集合框架中的泛型}

  \begin{itemize}[<+-|alert@+>]
  \item 泛型允许编译器实施由开发者设定的附加类型约束，将类型检查
    从{\hei\Blue 运行时挪到编译时}进行，这样类型错误就可以在编译时暴露
    出来，而不是在运行时才发作（抛出ClassCastException运行异常）。
  \item 创建集合容器时规定其允许保存的元素类型，然后由编译器负责添加元
    素的类型合法性检查，在取用集合元素时则不必再进行造型处理。
  \end{itemize}
\end{frame}

\begin{frame}[fragile]
  \frametitle{集合框架中的泛型}

  \tta{在Vector中使用泛型}

  \codeset{sample.generics.VectorGenericsSample.java}
  
  \tta{在Hashtable中使用泛型}

  \codeset{sample.generics.HashtableGenericsSample.java}
  
\end{frame}

\begin{frame}[fragile] % [fragile]参数使得能够插入代码
  \frametitle{泛型的向后兼容性}

  \begin{block}{}
    \begin{itemize}
    \item Java语言中的泛型是维护向后兼容的，完全可以不采用泛型、而继续
      沿用过去的做法。
    \item 这些未加改造的旧式代码将无法使用泛型带来的便利和安全性。
    \end{itemize}
  \end{block}


  
  未启用泛型机制的代码在高版本编译器中会输出如下形式的编译提示信息：
  
  \begin{stdoutCode}
    注: VectorGenericsSample.java使用了未经检查或不安全的操作。
    注: 有关详细信息, 请使用 -Xlint:unchecked 重新编译。
  \end{stdoutCode}

  \pause
  
  可以使用SuppressWarnings注解关闭编译提示或警告信息：
  
  \begin{javaCode}
    @SuppressWarnings({"unchecked"})  
    public class VectorGenericsSample {
      ...
    }
  \end{javaCode}
\end{frame}




\section{泛型类与泛型方法}

\begin{frame}[fragile] % [fragile]参数使得能够插入代码
  \frametitle{类型参数}

  \begin{alertblock}{类型参数和泛型类}
    形如Vector<String>，其中，尖括号括起来的部分被称为{\Red\hei 类型参
      数}，而这种由类型参数修饰的类型则被称为{\Blue\hei 泛型类}。
  \end{alertblock}

  \tta{注意}

  应在声明泛型类变量和创建对象时均给出类型参数，且两者应保持一致。
\end{frame}

\begin{frame}[fragile] % [fragile]参数使得能够插入代码
  \frametitle{类型参数}

  使用类型参数$E$进行泛型化处理的java.util.Vector类的定义代码摘要如下：
  
  \begin{javaCode}
    public class Vector<E> --- {
      public void addElement(E obj) { --- }
      public E elementAt(int index) { --- }
    }  
  \end{javaCode}

  \pause
  这里的$E$也称为“形式类型”参数。在实际使用该泛型类时，我们需要指定相应的
  具体类型，即实际类型参数。
  
  \begin{javaCode}
    Vector<String> v = new Vector<String>();  
  \end{javaCode}

  \pause
  编译器遇到Vector<String>类型变量时，即知道此Vector变量/对象的类型参
  数$E$已经被绑定为String类型，进而也就确定其addElement()方法的参数
  和elementAt()方法的返回值均为String类型。
\end{frame}




\begin{frame}[fragile] % [fragile]参数使得能够插入代码
  \frametitle{定义泛型类}

  \codeset{package sample.generics.userdefined}

  {\kai 代码示例中的泛型类PersonG可以在使用时通过类型参数T指定其属性secrecy的具体类型（以及该
    属性相应存/取方法的参数和返回值类型），进而提供了通用的信息存储能力。}

  \pause

  \tta{形式类型参数的编程惯例}

  \begin{table}
    \footnotesize
    \setlength{\extrarowheight}{1.2mm}
    \begin{tabular}{c|p{6cm}}
      {\bf 符号} & {\bf 意义}  \\
      K & 键，比如映射的键\\
      V & 值，比如List和Set的内容，或者Map中的值\\
      E & 元素，比如Vector<E>\\
      T & 泛型\\
    \end{tabular}
  \end{table}
\end{frame}

\begin{frame}[fragile] % [fragile]参数使得能够插入代码
  \frametitle{定义泛型类}

  \tta{使用受限制的类型参数}

  泛型机制允许开发者对类型参数进行附加约束。

  \begin{javaCode}
    import java.utils.Number;

    public class Point<T extends Number> {
      private T x;
      private T y;
      public Point() {}
      public Point(T x, T y) {
        this.x = x;
        this.y = y;
      }
      ...
    }
  \end{javaCode}

  \pause

  类型参数不能为基本数据类型，而java.lang.Number是所有数值型封装类（
  如Integer、Float、Double等）的父类型，于是限制泛型类Point的类型参数必
  须为Number或其子类类型，并使用extends关键字来标明这种继承层次上限制。
\end{frame}

\begin{frame}[fragile] % [fragile]参数使得能够插入代码
\frametitle{定义泛型方法}

与泛型类的情况类似，{\Red\hei 方法也可以被泛型化，且无论其所属的类是否为泛型类。}

\samp{泛型方法示例}

\begin{javaCode}
  public class Tool {
    public <T>T evaluate(T a, T b) {
      if(a.equals(b)) 
      return a;
      else
      return null;
    }
  }  
\end{javaCode}

\pause
\begin{itemize}
\item 上述代码中方法evaluate()声明中的“<T>”用于标明这是一个{\hei 泛型方
    法}。
\item 类型T是可变的，不必显示的告知编译器T具体取何值，但出现多处（两个
  形参、一个返回值类型）的这些值必须都相同。
\end{itemize}
\end{frame}

\begin{frame}[fragile] % [fragile]参数使得能够插入代码
\frametitle{定义泛型方法}

\tta{使用泛型方法，而不是定义泛型类的原则}

\begin{enumerate}[<+-|alert@+>]
\item 不涉及到类中的其他方法时，则应将之定义为泛型方法，因为泛型方法的
  类型参数是局部性的，这样可以简化其所在类型的声明和处理开销。
\item 要施加类型约束的方法为静态方法时，只能将之定义为泛型方法，因为静
  态方法不能使用其所在类的类型参数。
\end{enumerate}
\end{frame} 

\begin{frame}[fragile] % [fragile]参数使得能够插入代码
  \frametitle{对泛型的理解}
  
  \begin{itemize}
  \item 泛型类可以理解为具有广泛适用性、尚未最终定型的类型。
  \item Person<String>和Person<Double>属于同一个类，但确是不同的类型。
  \item 同一个泛型类与不同的类型参数复合而成的类型间并不存在继承关系，即使是类型参数间存在
    继承关系时也是如此。\\
    {\kai 如：Vector<String>不是Vector<Object>的子类。}
  \end{itemize}
\end{frame}

\section{处理泛型类型}

\begin{frame}[fragile] % [fragile]参数使得能够插入代码
  \frametitle{遍历泛型Vector集合}

  \tta{遍历Vector<String>类型集合}
  
  \begin{javaCode}
    public void overview(Vector<String> v) {
      for(String o: v) {
        String.out.println(o);
      }
    }
  \end{javaCode}

  \pause

  \notice{遍历其他类型参数的Vector集合元素该怎么办？}

  难道要定义overview(Vector<Person> v)、overview(Vector<Integer> v) ...
  
  显然过于繁琐，引用泛型机制后代码的通用性似乎不如从前？
\end{frame}

\begin{frame}[fragile] % [fragile]参数使得能够插入代码
  \frametitle{泛型类型的处理方法}

  \tta{可能的处理方法（不要使用）}

  将遍历方法的形参定义为不带任何类型参数的原型类型Vector，但这样会破坏已有的类型安全性。

  \begin{javaCode}
    public void overview(Vector v) {
      for(Object o: v) {
        String.out.println(o);
      }
    }
  \end{javaCode}

  \pause
  
  \tta{使用通配符}

  为了解决类似泛型遍历的问题，Java泛型机制中引入了通配符“？”。

  \begin{javaCode}
    public void overview(Vector<?> v) {
      for(Object o : v) {
        System.out.println(o);
      }
    }
  \end{javaCode}
\end{frame}

\begin{frame}[fragile] % [fragile]参数使得能够插入代码
\frametitle{类型通配符}
\wxd{使用类型通配符的好处}

\begin{enumerate}
\item Vector<?>是任何泛型Vector的父类型，因此可以将Vector<String>、Vector<Integer>、
  Vector<Object>等作为实参传给overview(Vector<?> v)方法处理；
\item Vector<?>类型的变量在调用方法时是受到限制的——凡是必须知道具体类型参数才能进行的操作
  均被禁止。
\end{enumerate}

\begin{javaCode}
Vector<String> vs = new Vector<String>();
vs.add("Tom");
Vector<?> v = vs;
v.add("Billy");  // 非法，编译器不知道具体类型参数
Object e = v.elementAt(0); // 合法，允许检索元素，此时读取的元素均当作 Object 类型处理
System.out.println(e);
\end{javaCode}
\end{frame}

\begin{frame}[fragile] % [fragile]参数使得能够插入代码
\frametitle{类型通配符}

上述限制不等同于将Vector<?>变为“只读”，在不需要编译器确定类型参数的情况下也是可以修改集
合内容的，例如：
\begin{javaCode}
Vector<String> vs = new Vector<String>();
vs.add("Tom");
vs.add("Billy");
Vector<?> v = vs;
v.remove(new Integer(200));  // 形参为 Object 类型，运行不受影响
v.clear(); // 不需要参数，运行不受影响
\end{javaCode}
\end{frame}

\begin{frame}[fragile]
  \frametitle{本节习题}

  \wxd{简答题}
  \begin{enumerate}
  \item 什么是泛型？泛型有什么作用？
  \end{enumerate}

  \wxd{小编程}
  \begin{enumerate}
  \item 自行定义一个泛型类，尝试写出代码。
  \item 练习泛型化集合类型基于通配符的遍历方法。
  \end{enumerate}

\end{frame}

% TKS %%%%%%%%%%%%%%%%%%%%%%%%%%%%%%%%%%%%%%%%%%%%%%
\begin{frame}[focus]
\centering
{\Huge {THE END}} \\
\vspace{5mm}
{\Large wangxiaodong@ouc.edu.cn} \\
\end{frame}
%%%%%%%%%%%%%%%%%%%%%%%%%%%%%%%%%%%%%%%%%%%%%%%%%%%%
\end{document}